\section{What the existing measurements support}
\label{sec:evidence}
The purpose of these measurements is to check the reasoning, not to fit every equation until
one method appears to win. The experiments have different scopes, so their numbers must not be
merged into one unqualified performance claim.

\Needspace{12\baselineskip}
\subsection{The one-million-passage scaling study}
The archived study used one flat cluster, a fixed 256-bit score, 100 binary candidates,
leaf size 32, deterministic branching, cached query embeddings, and one CPU thread.
There were 100 validation queries for selecting configurations and 200 held-out test queries,
with three timing repeats. It covered pools from 1,000 to one million real MS MARCO passages.

\begin{center}
\begin{tabular}{@{}lrrr@{}}\toprule
One-million-document test & Median latency & Binary recall & Fully scored rows\\\midrule
Packed scan & 25.14 ms & 100\% & 1,000,000\\
Branch, budget 32,768 & 244.98 ms & 99.26\% & about 163,766\\
Unlimited branch & 1,068.84 ms & 100\% & about 999,961\\\bottomrule
\end{tabular}
\end{center}
The branch reduced scoring at the finite budget, but did not win at matched target recall.
The whole study found no median or p95 crossover among the tested 95\% and 99\% target settings.
This conclusion belongs to that implementation, dataset, hardware and grid. It is not a theorem
that every bitplane representation loses.

The original implementation uses dense active masks for every node. Its measured split-word
counter omits leaf reads, which must be derived from the visited-node counts and cluster widths.
The physical mask size at one million rows is $15{,}625\times8=125{,}000$ bytes. A frontier with
20,000 such masks would contain 2.5 GB of masks alone. This example explains how query RAM can
be much larger than the packed-code payload; it is not a claimed observed frontier count.

The source summary is
\path{results/scaling-msmarco-1m/analysis/settings.csv}. Full run metadata and the validation
protocol live beside it. The report's verification script checks the numerical rows quoted here.

\Needspace{12\baselineskip}
\subsection{The FiQA study asked a broader pipeline question}
The denser FiQA run used 57,638 documents and 648 queries, with 583 settings including random
seeds. Its timings include routing, binary candidate search and float reranking, with embeddings
cached. Its quality measures include float Recall@10 and nDCG@10, rather than the scaling study's
binary Recall@100.

\begin{center}\small
\begin{tabular}{@{}lrrrr@{}}\toprule
Method & Probes & Node budget & Median ms & Float Recall@10\\\midrule
Sign-route binary scan & 16 & -- & 0.7524 & 0.6784\\
Sign-route branch & 32 & 2,048 & 0.6256 & 0.7006\\
Sign-route branch & 64 & 2,048 & 0.6466 & 0.7062\\
Float IVF reference & 8 & -- & 0.1661 & 0.9127\\\bottomrule
\end{tabular}\normalsize\end{center}
These selected examples show that probing more clusters with limited work inside them can be
a useful combined tradeoff. They do not isolate traversal efficiency: the selected pools differ.
Within the same selected buckets, unlimited branching was slower than scanning. The float IVF
row uses a different scorer and was a strong practical reference in that experiment.

Random exploration at probability 0.1 did not show a consistent improvement. Its purpose is to
change visit order; it does not provide a mathematical guarantee of higher recall. The full
FiQA report retains losing configurations, empty returns, and the relevant measurement scope.
Its source is \path{results/fiqa-2026-09-09-dense/README.md} and the accompanying data.

\Needspace{12\baselineskip}
\subsection{A new controlled check of the 24-bit prefix}
The small quality study used the same full weighted-binary score and row-ID tie rule for all
methods. There was no outer routing: $C=P=1$. Two synthetic datasets had 20,000 random binary
documents. The real-data case used the first 100,000 rows of the cached Nomic/MS MARCO sample.
All cases used $d=256$ and $h=24$.

One synthetic query distribution gave all coordinates equal variance. The other multiplied the
first 24 coordinates so that they held 80\% of \emph{expected} full-score variance. It was a
controlled favorable case for a prefix, not a model of what Matryoshka always does.

For each case, 16 queries estimated a score correlation and 16 different queries checked the
normal-model prediction. The native full scan agreed with an independent exact score sort on
all 96 query references. The short-key candidate set included all documents in keys tied at
the final prefix threshold, then used the exact full score to keep the top 100.

\begin{center}\small
\begin{tabular}{@{}lrrrr@{}}\toprule
Data, candidate target & Measured recall & Normal prediction & Gap & Median keys\\\midrule
Random equal variance, 1,000 & 23.00\% & 21.84\% & $-1.16$ pp & 846,156\\
Strong prefix, 1,000 & 95.00\% & 97.01\% & $+2.01$ pp & 846,156\\
Cached Nomic, 1,000 & 25.25\% & 17.41\% & $-7.84$ pp & 6,976\\
Cached Nomic, 5,000 & 52.13\% & 42.57\% & $-9.56$ pp & 35,504\\
Cached Nomic, 10,000 & 66.19\% & 58.46\% & $-7.73$ pp & 78,370\\\bottomrule
\end{tabular}\normalsize\end{center}
``Gap'' is prediction minus measurement in percentage points. Values are rounded from the
per-query evidence. The identical synthetic key counts are expected: the strong-prefix case
scales all first-24 query weights by the same factor, which changes neither their key order nor
the prefix candidate sets. It changes their importance in the \emph{full} score.

These results support two different conclusions. First, a short prefix is not inherently
inaccurate: the constructed favorable case reaches 95\% recall with 1,000 candidates.
Second, a simple statistical prediction can still miss real recall by several points. An
error of that size is unacceptable as a certificate of a 99\% recall target.

\fig{recall-validation}{Measured short-key recall versus the checked normal approximation. Each panel has its own explicitly stated document population.}{1}

The key counts were calculated from the actual query weights by meet-in-the-middle counting,
including empty keys. The candidate set was constructed using an exhaustive prefix-score
calculation for quality evaluation. \textbf{That computation is not an implemented hash-search
latency benchmark.} Its wall time must not be plotted as the latency of Approach C.

In the same favorable synthetic case, native B at budget 512 recovered about 99.38\% of the
binary top 100 while fully scoring about 4,122 documents on average. In the real-data case,
B reached about 58.06\%, 86.56\% and 99.56\% at budgets 512, 2,048 and 8,192. These are quality
and work observations for a flat pool with a 1,000-candidate heap. They do not establish a
latency ranking between an implemented native B and a hypothetical C.

\Needspace{12\baselineskip}
\subsection{Check the new derivation against independently sampled corpora}
\label{sec:recall-checks}
The small arithmetic checker enumerates ordered corpora and sorts their documents directly.
This includes duplicate codes, score ties, empty prefixes, full prefixes, $K=1$, and $K=N$.
All twenty settings agree with Equation~\eqref{eq:finite-top-k} to numerical precision. The
16-bit worked example is then calculated from the formula rather than by enumerating every
possible corpus.

A second check uses eight-bit documents, $N=64$, $K=5$, and the first three bits as lookup
coordinates. A document passes when at most one of those three bits disagrees with the query.
For each distribution, the script calculates the exact expected recall by grouping all possible
patterns by full penalty, then samples 10,000 independent corpora and checks the returned top
five directly. The query has positive coordinates with the following magnitudes:
\begin{itemize}
\item Equal weights: $[1,1,1,1,1,1,1,1]$; all bits independent and balanced.
\item Strong prefix: $[4,3,2,1,1,1,1,1]$; the same independent document distribution.
\item Weak prefix: $[1,1,1,4,3,2,2,2]$; the same independent document distribution.
\item Paired bits: all query magnitudes one, with the last four document bits repeating the
first four. Only sixteen patterns can occur.
\end{itemize}
These are deliberately simple assumptions. They let us change one cause at a time without
calling the resulting distribution a model of a particular embedding system.

\begin{center}\small
\begin{tabular}{@{}lrrr@{}}
\toprule
Case & Exact stated model & Sampled mean & Assume independent bits\\\midrule
Equal weights & 93.847\% & 93.784\% & 93.847\%\\
Strong prefix & 99.999\% & 100.000\% & 99.999\%\\
Weak prefix & 74.248\% & 74.430\% & 74.248\%\\
Paired bits & $\approx100\%$ & 100.000\% & 93.847\%\\\bottomrule
\end{tabular}\normalsize
\end{center}
The candidate rule is identical in every row. It returns an expected 32 of the 64 documents,
but its recall changes with the score weights and document distribution. The last row isolates
the failed assumption: individual bits remain balanced, yet assuming they are independent
underestimates recall by about 6.15 percentage points.

\fig{recall-model-checks}{Exact calculations and independently sampled corpora agree for the stated distributions. The paired-bit row shows the error from using the wrong distribution.}{1}

Each sampled corpus contributes one recall value in $[0,1]$. A conservative uncertainty bound
for their mean follows from
\[
\Pr(|\overline r-\E[r]|\ge\epsilon)\le2e^{-2n\epsilon^2}.
\]
Setting the right side to $0.05$ gives
$\epsilon=\sqrt{\log(40)/(2n)}$, about 1.36 percentage points at $n=10{,}000$.
The plotted intervals use this bound and are clipped to $[0,1]$. This remains meaningful when
no misses were observed. A zero sample variance by itself would produce a misleading zero-width
normal interval near 100\%. The bound assumes independent sampled corpora; it does not turn
ten thousand simulations into evidence about real embeddings.

For the first twelve sampled corpora in each case, the check also calls the current native
scan and bitplane search. All 48 full scans and all 48 unlimited branch searches match the
independently sorted top five. With a 16-node budget and leaf size four, mean branch recalls
are 80.00\%, 98.33\%, 100.00\%, and 85.00\% in the table's order. At exploration probability
0.1 and the same budget, they are 75.00\%, 98.33\%, 100.00\%, and 85.00\%. These small examples
show that the budget and weights affect which neighbors survive. They are not estimates of
performance on a large dataset or evidence that random exploration always hurts.

The script separately reaches the same fixed prefix-radius candidate set through key grouping
and through the union of bitplane leaves. All 48 candidate sets agree exactly. This checks the
shared-set argument; it does not claim that budgeted best-first traversal visits those same
leaves. No latency ranking is inferred from these quality checks.

\Needspace{12\baselineskip}
\subsection{Put an uncertainty interval around the saved prediction errors}
The earlier study already showed that the normal approximation could miss real recall.
To see how variable that gap is across queries, resample its sixteen check queries with
replacement 10,000 times. For each resample, average the paired difference
$\widehat r(q)-r(q)$. The middle 95\% of those averages gives the bootstrap interval below.
A percentage point is an absolute difference: predicting 30\% when the observation is 25\%
is an error of $+5$ points.

\begin{center}\small
\begin{tabular}{@{}lrrr@{}}
\toprule
Case & Candidate target & Mean error (points) & 95\% interval\\\midrule
Random signs & 1,000 & $-1.16$ & $[-3.72,\ 1.41]$\\
Strong prefix & 1,000 & $+2.01$ & $[0.08,\ 4.14]$\\
Cached Nomic & 1,000 & $-7.84$ & $[-14.47,\ -1.54]$\\
Cached Nomic & 5,000 & $-9.56$ & $[-18.25,\ -0.81]$\\
Cached Nomic & 10,000 & $-7.73$ & $[-16.17,\ 0.64]$\\\bottomrule
\end{tabular}\normalsize
\end{center}
This is a reanalysis of previously examined queries, not a fresh test set. The fitted score
correlation is held fixed, so the intervals do not include uncertainty from estimating that
correlation. They also do not cover corpus changes, model changes, or selection among many
configurations. When every checked query gives the same outcome, bootstrap variation can
vanish; that is not proof that future queries have no variation.

A proposed requirement of at most two points of mean error fails on every cached-Nomic setting
shown here. Even the random-sign row, whose observed mean error is small, has an interval that
extends beyond that tolerance. We should therefore retain this normal model as an explanation
of a possible mechanism, rather than use it to certify real 99\% recall.

Improving the prediction requires new evidence about the relevant score distribution. For
example, estimate the joint full and prefix penalties on calibration queries, freeze the
candidate policy, and check it on separate queries. If an empirical distribution is used, its
resolution in the nearest-score tail must be reported. A small sample can miss rare but
important patterns; moving from 100,000 to a billion documents makes that issue more severe.
The current extension does not claim a newly validated real-data predictor.

\Needspace{12\baselineskip}
\subsection{What was not validated}
The study does not validate:
\begin{itemize}
\item real recall at a billion documents;
\item a routing-recall model across arbitrary cluster counts;
\item a normal tail for every query or a fixed percentage of Matryoshka prefix information;
\item a concrete hash enumerator's timing, peak queue size, or SIMD speed;
\item float-vector recall, human relevance, or final answer accuracy from binary recall alone.
\end{itemize}
Float-vector top-100 results and binary-score top-100 results answer different questions.
This report keeps their comparisons separate. A method that searches binary scores exactly
can still miss neighbors under the original float-vector score.

\Needspace{12\baselineskip}
\subsection{How to extend the checks fairly}
Freeze the representation, query set, filtering task and reference score first. Tune on one
query set; check selected configurations on a different one. For every method, retain failures
and empty returns, and record counts as well as times. For C, implement the actual enumerator
before claiming its measured latency. For B, record split words, leaf words, documents scored
and peak frontier memory separately.

At each target recall, compare measured latency quantiles using configurations that actually
reach that target. Use the same query population and report the uncertainty of quality estimates.
Query-level resampling is appropriate; counting repeated timings as independent quality samples
is not. A few measured cases are useful evidence about mechanisms, not a complete characterization
of every data distribution.
